\documentclass[11pt]{article}
\usepackage[margin=1in]{geometry}
\usepackage{amsmath,amssymb,amsthm,hyperref}
\setlength{\emergencystretch}{2em}
\newtheorem{theorem}{Theorem}
\newtheorem{lemma}{Lemma}
\title{A deterministic positive-weight counterexample\\to Conjecture 00000009876}
\author{ziangni-sys}
\date{October 8, 2026}
\begin{document}
\maketitle

\paragraph{Claim and scope.}
The English statement asserts that passage-time variances of general
finite-variance-weight first-passage percolation grow as the one-half power
of distance, and that the exponent $1/2$ is universal under finite
exponential moments. The Chinese statement makes the same universal
assertions. Neither version requires a nonconstant edge-weight law.
We disprove both assertions as written using strictly positive,
independent, identically distributed weights on the standard square lattice.
The example does not address a modified conjecture that explicitly excludes
deterministic edge laws. A false universal clause suffices to disprove the
conjunction; no claim concerning the heavy-tail clause is needed.

\paragraph{The probability model.}
Let $V=\mathbb Z^2$, and join two vertices when they differ by one in exactly
one coordinate. Let $E$ be the set of its undirected edges. On the one-point
probability space $(\{*\},\delta_*)$, set
\[
             W_e(*)=1 \qquad(e\in E).
\]
Every $W_e$ has law $\delta_1$, is strictly positive, has finite second moment,
and satisfies $\mathbb E e^{tW_e}=e^t<\infty$ for every real $t$.
These variables are independent, including in the full finite-family sense:
for any finite $F\subset E$ and Borel sets $B_e$,
\[
 \mathbb P\left(\bigcap_{e\in F}\{W_e\in B_e\}\right)
 =\prod_{e\in F}\mathbf1_{\{1\in B_e\}}
 =\prod_{e\in F}\mathbb P(W_e\in B_e).
\]
Thus this is an iid first-passage-percolation model, despite its degenerate
distribution. In particular, neither lack of positivity nor lack of moments
causes the counterexample.

\paragraph{Actual passage times.}
For a finite nearest-neighbor walk $p=(v_0,\ldots,v_k)$, its cost is
$C(p)=\sum_{j=0}^{k-1}W_{\{v_j,v_{j+1}\}}=k$.
Define passage time by the usual infimum over all such walks:
\[
 T(x,y)=\inf\{C(p):p\text{ is a finite nearest-neighbor walk from }x\text{ to }y\}.
\]
Allowing walks with repeated vertices makes no difference here: all costs
are positive. Put $T_n=T((n,0),(0,0))$, where $n\geq0$ is an integer.
Undirected symmetry also gives the same passage time in the reverse direction.

\begin{lemma}
For every $n\geq0$, $T_n=n$ at every outcome.
\end{lemma}
\begin{proof}
Each edge can decrease the first coordinate by at most one. Thus a walk of
length $k$ from $(n,0)$ to $(0,0)$ must have $n\leq k$. More explicitly,
the inequality $v_{j,1}\leq v_{j+1,1}+1$ sums to $n\leq k$.
Consequently every admissible cost is at least $n$. The horizontal walk
$(n,0),(n-1,0),\ldots,(0,0)$ has exactly $n$ edges and cost $n$.
The set in the infimum is therefore nonempty and bounded below, its infimum
equals $n$, and the infimum is attained.
\end{proof}

\begin{theorem}
The passage-time variances are identically zero. In particular there are no
$c>0$ and $N\in\mathbb N$ such that
\[
       c\sqrt n\leq\operatorname{Var}(T_n)\quad\hbox{for all }n\geq N.
\]
Hence the universal one-half growth assertion is false.
\end{theorem}
\begin{proof}
The preceding lemma gives $\mathbb E T_n=n$ and
$\operatorname{Var}(T_n)=\mathbb E[(T_n-n)^2]=0$.
If the displayed inequality held, substituting $n=N+1$ would give a strictly
positive left side and zero right side, a contradiction.
\end{proof}

\paragraph{Meaning of ``grow as''.}
The contradiction applies to $\operatorname{Var}(T_n)\sim c\sqrt n$ with
$c>0$, to two-sided comparability with $\sqrt n$, and to any assertion of
a positive growth exponent $1/2$. Zero variance cannot have that exponent.
An upper bound $O(\sqrt n)$ alone would be consistent with this example,
but would not assert growth as a one-half power or universality of that
exponent, as the source does.

\paragraph{Formal verification.}
The accompanying Lean 4 project defines the square lattice as a
\texttt{SimpleGraph} on $\mathbb Z\times\mathbb Z$, its actual undirected
edge type, the Dirac probability measure, and the edge random variables.
It proves their joint independence with Mathlib's \texttt{iIndepFun}, their
common law, positivity, second-moment integrability, and exponential-moment
integrability for every real parameter. Costs are recursively summed along
\texttt{SimpleGraph.Walk}; passage times use the real infimum of all walk
costs. The coordinate lower bound and explicit horizontal walk establish
\texttt{passage\_eq\_distance}. Mathlib's measure-theoretic variance yields
\texttt{passage\_variance\_zero}, followed by the explicit negation
\texttt{no\_positive\_sqrt\_lower\_bound}. No asymptotic conclusion, shortest-path
certificate, or variance identity is assumed as a hypothesis.
\end{document}
